\section{Conclusion}

We set out to detect spurious features by their emergence uniformity --- and discovered why this elegant idea fails on pretrained models.

\subsection{Summary}

CV of linear probe accuracy trajectories on frozen CLIP features shows no discriminative power for spurious vs.\ core features (AUC = 0.0). The root cause is feature saturation: CLIP has already learned both concepts.

\subsection{Key Insight}

Emergence uniformity is a training-time phenomenon that cannot be recovered from frozen pretrained representations. You cannot measure how fast someone learned by looking at their final exam score.

\subsection{Contributions}

\begin{enumerate}
    \item \textbf{Negative result}: CV on frozen CLIP features does not distinguish spurious from core features
    \item \textbf{Mechanistic explanation}: Feature saturation eliminates emergence dynamics
    \item \textbf{Design principle}: Emergence-based detection requires training-time measurement
\end{enumerate}

\subsection{Future Directions}

\begin{enumerate}
    \item Training-time CV measurement
    \item Intermediate layer probing
    \item Loss-based emergence signals
    \item Alternative feature extractors
\end{enumerate}

If emergence uniformity proves measurable during training, it could enable single-run, annotation-free spurious correlation mitigation --- the original goal of EUR.
